% !TEX program = lualatex
%
%
% Generated by new_lecture_note.sh
\documentclass[11pt]{article}

\usepackage{fontspec}
\usepackage{microtype}
\usepackage{mathtools,amssymb,amsfonts,amsthm,bm}
\usepackage{enumitem}
\usepackage{booktabs,array,xcolor}
\usepackage{geometry,fancyhdr,setspace}
\usepackage[hidelinks]{hyperref}
\usepackage[nameinlink,noabbrev]{cleveref}

\geometry{margin=1in,headheight=15pt}
\setstretch{1.12}
\setlength{\parindent}{0pt}
\setlength{\parskip}{0.55em plus 0.12em minus 0.08em}
\setlength{\abovedisplayskip}{0.8em plus 0.2em minus 0.1em}
\setlength{\belowdisplayskip}{0.8em plus 0.2em minus 0.1em}
\allowdisplaybreaks[3]
\setlist{leftmargin=*,topsep=0.35em,itemsep=0.15em,parsep=0pt}

% ---------- Metadata ----------
\newcommand{\studentname}{Keshav Ramamurthy}
\newcommand{\coursename}{Math 113}
\newcommand{\lecturenumber}{02}
\newcommand{\lecturetitle}{Lecture 02}
\newcommand{\lecturedate}{\today}

% ---------- Reusable notation ----------
\newcommand{\N}{\mathbb{N}}
\newcommand{\Z}{\mathbb{Z}}
\newcommand{\Q}{\mathbb{Q}}
\newcommand{\R}{\mathbb{R}}
\newcommand{\C}{\mathbb{C}}
\newcommand{\F}{\mathbb{F}}
\newcommand{\K}{\mathbb{K}}
\newcommand{\E}{\mathbb{E}}
\newcommand{\Prob}{\mathbb{P}}
\newcommand{\1}{\mathbf{1}}
\newcommand{\eps}{\varepsilon}
\newcommand{\dd}{\,\mathrm{d}}
\newcommand{\abs}[1]{\left\lvert#1\right\rvert}
\newcommand{\norm}[1]{\left\lVert#1\right\rVert}
\newcommand{\ip}[2]{\left\langle#1,#2\right\rangle}
\newcommand{\set}[1]{\left\{#1\right\}}
\newcommand{\given}{\,\middle\vert\,}
\newcommand{\indep}{\perp\!\!\!\perp}
\newcommand{\lto}{\longrightarrow}
\newcommand{\deriv}[2]{\frac{\mathrm{d}#1}{\mathrm{d}#2}}
\newcommand{\pderiv}[2]{\frac{\partial#1}{\partial#2}}
\DeclareMathOperator{\Aut}{Aut}
\DeclareMathOperator{\Hom}{Hom}
\DeclareMathOperator{\Ker}{ker}
\DeclareMathOperator{\im}{im}
\DeclareMathOperator{\rank}{rank}
\DeclareMathOperator{\nullity}{nullity}
\DeclareMathOperator{\Span}{span}
\DeclareMathOperator{\Var}{Var}
\DeclareMathOperator{\Cov}{Cov}

% ---------- Note environments ----------
\newtheoremstyle{notestyle}{0.7em}{0.7em}{\itshape}{}{}{\bfseries}{.5em}{}
\theoremstyle{notestyle}
\newtheorem{theorem}{Theorem}[section]
\newtheorem{lemma}[theorem]{Lemma}
\newtheorem{proposition}[theorem]{Proposition}
\newtheorem{corollary}[theorem]{Corollary}
\theoremstyle{definition}
\newtheorem{definition}[theorem]{Definition}
\newtheorem{example}[theorem]{Example}
\newtheorem{remark}[theorem]{Remark}
\newenvironment{claim}{\par\noindent\textbf{Claim.}\ }{\par}
\newenvironment{idea}{\par\noindent\textit{Idea.}\ }{\par}
\newenvironment{proofsketch}{\par\noindent\textbf{Proof sketch.}\ }{\par}

% ---------- Header ----------
\pagestyle{fancy}
\fancyhf{}
\fancyhead[L]{\coursename}
\fancyhead[C]{\lecturetitle}
\fancyhead[R]{\studentname}
\fancyfoot[C]{\thepage}

\title{\vspace{-1.5em}\textbf{\coursename -- \lecturetitle}}
\author{\studentname}
\date{\lecturedate}

\begin{document}
\maketitle

\section{Groups}
% Binary Operator with additional axioms of binary structures.
A group is a binary algebraic structure $(G, *)$ such that 

\begin{itemize}
    \item $*$ is associative $\forall a, b, c \in G, a * (b * c) = (a * b) * c$
    \item $\exists e \in G$ such that e is an identity element $e * a = a * e = a$
    \item $\forall a \in G,  \exists a^{-1} \in G$ (each element has an inverse)
    \item In this class, assume it's a two sided inverse, namely $a a^{-1} = a^{-1} a = e$
\end{itemize}
\begin{definition}[Abelian Group]
\label{def:label}
A group $G$ is abelian if $*$ is commutative ($\forall a, b \in G,  a * b = b * a.$)
\end{definition}
% In this class, assume it's a two sided inverse, namely $a a^{-1} = a^{-1} a = e$
% \noindent\textbf{Proof.} 
% \noindent\textbf{Proof.} 
\section{Definitions and notation}
\begin{theorem}[Cancellation Law for Groups]
\label{thm:label}
For a group $(G, *)$ and $a, b, c \in G$, we have $$a * b = a * c \iff b = c$$
\end{theorem}
\begin{proof}
\begin{align*}
    a * b = a * c 
    a^{-1} * (a * b) &= a^{-1} * (a * c)\\
 (a^{-1} * a) * b &= (a^{-1} * a) * c\\ 
 e * b &= e * c \rightarrow b = c\\ 
\end{align*}
\end{proof}

\begin{theorem}[Uniqueness]
\label{thm:label}
For a group $(G, *)$ and $a, b \in G$, then $$a * x = b \text{ has a unique solution } x\in G$$
\end{theorem}
\begin{proof}
Existence: x = $a^{-1}b$ is a solution.
\par Uniqueness: Suppose for sake of contradiction, $\exists y \in G$ such that $a * x =
b = a * y$ \par However, by the cancellation law $a * x = a * y \iff x = y$ a contradiction.
\end{proof}
\begin{theorem}[Uniqueness of Inverses]
\label{thm:label}
For $a \in G$ of a group $(G,*)$, it has a unique inverse $a^{-1}.$
\end{theorem}
\begin{proof}
Any inverse is a solution to $a * x = e$, and by the uniqueness theorem, it has a unique
solution
\end{proof}
\begin{corollary}[Corollary to Previous Theorem]
\label{cor:label}
For $a, b \in G$ of a group $(G, *)$, we have $$(a * b)^{-1} = b^{-1} a^{-1}$$
\end{corollary}
\begin{proof}
We check it: 
\begin{align*}
    (a * b) * (b^{-1} * a^{-1})
        &= a * (b * b^{-1}) * a^{-1} \\
        &= a * e * a^{-1} \\
    &= a * a^{-1} \\
    &= e\\
\end{align*}

% efaehla ia salih asli fas e$5 * 3 = 15$
\end{proof}
\begin{definition}[First definition]
Write the definition and one sentence explaining why it matters.
\end{definition}

\section{Claims, theorems, and proof sketches}
\begin{claim}
State the claim in one precise sentence.
\end{claim}

\begin{proofsketch}
Record the key idea, the first useful equation, and the step that still needs checking.
\end{proofsketch}

\section{Examples and calculations}
\begin{example}[Lecture example]
Write the setup, the calculation, and the conclusion.
\end{example}

\section{Questions and follow-up}
\begin{itemize}
  \item What is still unclear?
  \item Which exercise or reading should I revisit?
  \item TODO: 
\end{itemize}

\end{document}
